\section{Arithmetic models for the original principal uniformization}
\label{sec:principal-uniformization-v48}

\status{This packet constructs arithmetic suppliers and individual
moduli points for the original principal-uniformization theorem.
The original field, host,
perfect alternating lattices, level ideals and Hermitian space are
retained.  The arithmetic classical input is Hasse--Minkowski,
already identified in MF4; the final subsection separately names
the classical complex-torus and algebraic-geometric foundations
used for individual moduli points.  The proofs below do not assume a Hermitian
local-to-global theorem or the desired principal uniformization.
The precise remaining holomorphic-family and normalized-frame inputs
are recorded at the end.  In particular, an arithmetic lattice model
alone does not establish the common principal-chart obligation.}

Put \(F=E^{\bar{\ } }\), \(R=\mathcal O_E\), and
\(R_0=\mathcal O_F\).  The original \texttt{Host} supplies a CM field
\(E\), an everywhere unramified quadratic extension \(E/F\), and
\(\mathrm{rt}^2=-d\), where \(d\) is a positive prime and
\(d\equiv7\pmod {16}\).  Its chosen complex square root is
\(\mathrm{rt}_0\).  Write
\[
 \Psi=\{\tau:E\longrightarrow\mathbb C:
                       \tau(\mathrm{rt})=\mathrm{rt}_0\}.
 \tag{PUN.1}
\]
Complex conjugation takes \(\mathrm{rt}\) to \(-\mathrm{rt}\), so
\(\Psi\) contains exactly one member of each conjugate pair.
Throughout this packet Hermitian forms are \emph{linear in the first
variable}, as in the original \texttt{herm}.  An original datum
\(D\) consists of a finite projective \(R\)-module \(L\), its perfect
alternating \(\mathbb Z\)-pairing \(B\), the adjoint identity
\(B(ax,y)=B(x,\bar a y)\), and the original compatible positive
complex structure with multiplicities \((p,q)\).

An original level is denoted \(l=(M,\mathfrak b,A)\).  In particular
\(M\ge3\) is odd, \(A>0\), and
\[
 \mathfrak b\bar{\mathfrak b}=MR,\qquad
 \mathfrak b+\bar{\mathfrak b}=R,\qquad
                  N=l.M'=\operatorname{lcm}(M,A).
 \tag{PUN.2}
\]
No freeness of \(L\) over \(R\) is imposed.  It is finite free as a
\(\mathbb Z\)-module.  The empty matrices and zero lattice have their
usual meanings in rank zero.

\subsection{The norm equation and the same Hermitian space}

% BEGIN PUN1
\paragraph{PUN1. The literal quadratic norm theorem.}
\label{claim:PUN1}
The classical Hasse--Minkowski theorem in MF4 implies the whole
original \texttt{Published.HasseNormSt}: for every number field
\(K\), every \(\delta,x\in K^\times\), if
\(x=a^2-\delta b^2\) is solvable at every real embedding and every
finite completion, it is solvable in \(K\).

\begin{proof}
If \(\delta=s^2\), take
\(a=(x+1)/2\), \(b=(x-1)/(2s)\).  Otherwise apply
Hasse--Minkowski to the nondegenerate ternary quadratic form
\[
                  Q(U,V,W)=U^2-\delta V^2-xW^2.
 \tag{PUN.3}
\]
At the specified completions there is an isotropic vector with
\(W=1\).  At a complex completion a square root of \(x\) gives
one as well.  A nonzero global isotropic vector cannot have \(W=0\):
then \(V\ne0\) would make \(\delta=(U/V)^2\), and \(V=0\)
would force all three coordinates to vanish.  Division by \(W\)
gives the required solution.  This preserves every field and every
place quantified in the original statement, without an auxiliary
extension or an omitted dyadic case.
\end{proof}
% END PUN1

% BEGIN PUN2
\paragraph{PUN2. Hermitian local-to-global isometry by a quadratic induction.}
\label{claim:PUN2}
Let \(E/F\) be a CM quadratic extension and let \(h,h'\) be
nondegenerate Hermitian spaces of the same finite dimension.
If they are isometric at every finite completion of \(F\) and at
every real embedding, they are isometric over \(E\).

\begin{proof}
Diagonalize \(h'\) and let \(a\in F^\times\) be its first
coefficient.  The quadratic form \(q(v)=h(v,v)\) on the underlying
\(F\)-space of \(h\) is nondegenerate: its polar form is
\(\operatorname{Tr}_{E/F}h(v,w)\), whose radical is zero by the
nondegeneracy of \(h\) and of the field trace.  The form
\(q(v)-a t^2\) is locally isotropic everywhere.  Hasse--Minkowski
supplies a nonzero global zero \((v,t)\).

If \(t\ne0\), then \(h(v/t,v/t)=a\).  If \(t=0\), choose
\(w\) with \(h(v,w)=1\) and replace \(w\) by
\(w-h(w,w)v/2\).  The resulting \(w_0\) has norm zero and
\(h(v,w_0)=1\).  Thus \(v+(a/2)w_0\) has norm \(a\).
In either case split off a global line of norm \(a\).

At each completion, two vectors of the same nonzero norm are
carried to one another by an isometry.  For a field extension this
is the explicit reflection calculation of MF3, with the arguments
interchanged to convert its conjugate-linear-first convention to
the present convention.  If its first denominator vanishes, the
intermediate vector \(-w\) makes the alternative denominator
\(2a\ne0\).  For a split quadratic algebra a Hermitian space is
a perfect pairing between two vector spaces; extend the specified
nonzero paired vectors to dual bases.  This gives the same assertion
without dividing by a possible zero divisor.  Consequently the two
orthogonal complements are locally isometric at every place.
Induction gives an isometry of the complements, and adjoining the
chosen lines gives the desired isometry.  The dimension-zero case
starts the induction.  This argument uses only the stated quadratic
theorem; no separate Landherr theorem is being assumed.
\end{proof}
% END PUN2

% BEGIN PUN3
\paragraph{PUN3. The signature forced by the original positive complex structure.}
\label{claim:PUN3}
Suppose \(\phi:L\hookrightarrow E^r\) has full span, \(h=h_G\)
is nondegenerate, \(\bar\xi=-\xi\ne0\), and
\[
              B(x,y)=\operatorname{Tr}_{E/\mathbb Q}
                              (\xi h(\phi x,\phi y)).
 \tag{PUN.4}
\]
Then \(r=p+q\).  At \(\tau\in\Psi\), if
\(\operatorname{Im}\tau\xi>0\), \(h_\tau\) has \(q\)
positive and \(p\) negative directions; if the imaginary part is
negative, these multiplicities are reversed.  In particular two
original data of the same \((p,q)\), represented using the same
\(\xi\), have identical archimedean Hermitian signatures.

\begin{proof}
Tensoring the injection with \(E\) gives an isomorphism, so the
rank is the original rank \(p+q\).  Under
\(E\otimes_{\mathbb Q}\mathbb R\simeq\prod_{\tau\in\Psi}\mathbb C\),
the pairing (PUN.4) becomes
\(\sum_\tau2\operatorname{Re}(\tau\xi\,h_\tau)\).
The original complex structure \(J\) commutes with \(E\), hence
with the idempotents and complex scalars in this product.  It
therefore acts by a complex-linear operator \(J_\tau\) on each
\(\mathbb C^r\), with \(J_\tau^2=-1\).

The two real trace identities in \texttt{Datum.dom}, first on
\(E\) and then by real-linear extension on
\(E\otimes\mathbb R\), give
\[
 \operatorname{tr}_{\mathbb C}J_\tau=i(p-q).
 \tag{PUN.5}
\]
Together with \(J_\tau^2=-1\), this means that its \(+i\) and
\(-i\) eigenspaces have dimensions \(p\) and \(q\).
The invariance of \(B\), tested also after multiplying one argument
by arbitrary complex scalars in the same product factor, gives
\(h_\tau(J_\tau u,J_\tau v)=h_\tau(u,v)\).  Thus these two
eigenspaces are orthogonal.  Write \(\tau\xi=is\) with
\(s\in\mathbb R^\times\).  On them the original positivity
\(B(Ju,u)>0\) says respectively
\(-2s h_\tau(u,u)>0\) and \(2s h_\tau(u,u)>0\).
The asserted signature follows.  This also proves the required
diagonal normal form with first \(q\) signs positive when \(s>0\):
choose orthogonal bases in the two definite eigenspaces and rescale
each vector by the positive square root of its absolute norm.
Conjugate embeddings have the same signatures by conjugating all
coordinates.
\end{proof}
% END PUN3

% BEGIN PUN4
\paragraph{PUN4. Local isometry from the perfect integral trace pairing.}
\label{claim:PUN4}
Under the original \texttt{Host.unram}, let two full lattices in
rank \(r\) Hermitian spaces satisfy (PUN.4), with the same nonzero
\(\xi\), and let both resulting alternating integral pairings be
perfect.  At every finite place of \(F\) their Hermitian spaces
are isometric, including the places above \(2\).

\begin{proof}
Let \(C=\xi^{-1}\mathfrak D_{E/\mathbb Q}^{-1}\).  The exact
trace-duality calculation of PCL2--PCL3 gives, for either lattice,
\[
          L=\{v:h(v,L)\subseteq C\}.
 \tag{PUN.6}
\]
This identity commutes with localization and completion: a finite
projective lattice and its dual are finite projective, and the
pairing matrix may be computed in a local basis.  Fix a finite
place \(v\) of \(F\), put \(A=\widehat{\mathcal O}_{F,v}\)
and \(B=\mathcal O_E\otimes_{\mathcal O_F}A\).  The original
unramified hypothesis means that \(B\) is either \(A\times A\)
or the integers of an unramified quadratic extension.  In either
case every conjugation-stable fractional ideal of \(B\) is
\(cB\) for some \(c\in F_v^\times\): in the field case choose
a power of a uniformizer of \(F_v\), and in the split case
conjugation stability makes the two valuations equal.

The ideal \(C\) is conjugation stable, since
\(\bar\xi=-\xi\) and the different is stable.  Choose the same
\(c\) for both lattices.  Equation (PUN.6) says exactly that
\(c^{-1}h\) is a unimodular Hermitian form on a free \(B\)-module
of rank \(r\).  MF2 diagonalizes this form over \(B\) and removes
each diagonal unit by a unit norm.  Its proof uses the surjective
residue trace and successive corrections, so it includes residue
characteristic \(2\).  In the split case dual bases for the two
factors give the same conclusion.  Both scaled forms are therefore
isometric to the same standard form.  Multiplication by \(c\)
gives the asserted local isometry of the original forms.
\end{proof}
% END PUN4

% BEGIN PUN5
\paragraph{PUN5. Every component lattice embeds in the prescribed Hermitian space.}
\label{claim:PUN5}
Let \(D,D'\) be original data for the same host, with
\(D'.p=D.p\) and \(D'.q=D.q\).  Suppose the prescribed
\((G,\phi,\xi)\) satisfies the full original Hermitian-coordinate
contract (PUN.4), including injectivity, full span, Hermitian
symmetry and \(\det G\ne0\).  There is an \(R\)-linear
injection \(\phi':D'.L\hookrightarrow E^r\) with full span and
\[
        B'(x,y)=\operatorname{Tr}_{E/\mathbb Q}
                      (\xi h_G(\phi'x,\phi'y)).
 \tag{PUN.7}
\]
The matrix \(G\) and scalar \(\xi\) are unchanged.

\begin{proof}
For \(\xi\ne0\), PCL1 supplies Hermitian coordinates
\((G',\psi)\) for \(D'\) using this same \(\xi\).  The ranks
agree by PUN3.  The two forms have the same infinite signatures
by PUN3 and are isometric at every finite place by PUN4.  PUN2
provides an \(E\)-linear isometry
\(g:(E^r,h_{G'})\longrightarrow(E^r,h_G)\).
Take \(\phi'=g\psi\).  It is injective, has full span, and
satisfies (PUN.7) by the isometry identity.

For the literal source's additional case \(\xi=0\), (PUN.4)
and perfection imply \(D.L=0\); full span then gives \(r=0\).
The original rank equality and the equal signatures give
\(D'.L=0\), since a finite projective module over a domain of
rank zero is zero.  The unique map into \(E^0\) proves the
assertion.  Thus this is the full same-space producer, with a
quadratic local-to-global proof replacing its longer norm-ideal
implementation.
\end{proof}
% END PUN5

\subsection{The actual principal group and all lattice-model fields}

% BEGIN PUN6
\paragraph{PUN6. The original half-level condition is the full congruence condition.}
\label{claim:PUN6}
Let \(g\) be an \(R\)-linear automorphism of the original \(L\)
preserving \(B\).  If
\((g-1)\bar{\mathfrak b}L\subseteq ML\), then
\((g-1)L\subseteq ML\).  Consequently the original group
\(\operatorname{princ}(l,G,D,\phi)\) is exactly
\[
 \{g\in U(G):g\Lambda=\Lambda,
                        (g-1)\Lambda\subseteq N\Lambda\},
 \qquad \Lambda=\phi L,\quad N=\operatorname{lcm}(M,A).
 \tag{PUN.8}
\]

\begin{proof}
Choose \(i\in\mathfrak b\) and \(j\in\bar{\mathfrak b}\)
with \(i+j=1\).  Put \(y=gx-x\).  The assumption gives
\(jy=Me_1\).  For any \(z\in L\), write
\(g(\bar i z)=\bar i z+Me_2\); this is permitted because
\(\bar i\in\bar{\mathfrak b}\).  The adjoint and isometry
identities give
\[
 B(iy,z)=B(gx-x,\bar i z)=-M B(gx,e_2).
 \tag{PUN.9}
\]
Thus \(z\mapsto B(iy,z)/M\) is an integral linear functional.
Perfection gives \(w\in L\) representing it, and injectivity of
the pairing gives \(iy=Mw\).  Adding the equations for \(iy\)
and \(jy\) proves \(y=M(w+e_1)\).

An element stabilizing \(\Lambda\) induces an actual automorphism
of \(L\), because \(\phi\) is an isomorphism onto its range.
Its Hermitian isometry preserves \(B\), so the preceding argument
applies to the literal half-level condition.  The additional
\(A\)-condition gives
\((g-1)\Lambda\subseteq M\Lambda\cap A\Lambda=N\Lambda\),
where the last equality follows in a \(\mathbb Z\)-basis of
\(\Lambda\).  The converse follows from \(M\mid N\),
\(A\mid N\).  This proves equality with the original group,
not just a finite-index inclusion.
\end{proof}
% END PUN6

% BEGIN PUN7
\paragraph{PUN7. Finite index with the source's coset order.}
\label{claim:PUN7}
The actual principal group has finite index in
\(U(G,\Lambda)=\{g\in U(G):g\Lambda=\Lambda\}\).
There is a finite set \(S\subseteq\operatorname{GL}_r(E)\)
such that every \(g\in U(G,\Lambda)\) has an \(s\in S\) with
\(s^{-1}g\in\operatorname{princ}(l,G,D,\phi)\).

\begin{proof}
The set \(\Lambda/N\Lambda\) is finite since \(N>0\) and
\(\Lambda\) is finite free over \(\mathbb Z\).  Reduction gives
a homomorphism from \(U(G,\Lambda)\) to the automorphism group
of this finite module.  Its kernel is exactly (PUN.8).  Choose
one representative of each element of the finite image.  Equal
reductions of \(g\) and its representative \(s\) give
\(s^{-1}g\) in that kernel, which is the precise order required
by the original \texttt{Gamma\_index} field.
\end{proof}
% END PUN7

% BEGIN PUN8
\paragraph{PUN8. Freeness on the original full flag domain.}
\label{claim:PUN8}
For every \(g\) in the original principal group, every
\(x\in\operatorname{GG}(\Psi,G,p_0)\), and every original
upper-triangular \(b\in\operatorname{Bor}(\Psi,r)\),
\(g\cdot x=x b\) implies \(g=1\).

\begin{proof}
At each \(\tau\in\Psi\), the successive column spaces of
\(x_\tau\) form a complete flag whose restrictions of
\(h_\tau\) are nondegenerate: their determinants are precisely
the nonzero leading minors in the definition of \(\operatorname{GG}\).
An isometry preserving this flag preserves the one-dimensional
orthogonal difference between each two consecutive terms.  On
each such nondegenerate line it acts by a scalar of absolute value
one.  It is therefore diagonalizable in a fixed basis with all
powers bounded.  The equality \(g x=x b\) ensures flag
preservation.  Conjugate embeddings give the same bounds; hence
the real-linear operators induced by all powers of \(g\) are
bounded on \(\Lambda\otimes\mathbb R\).

Choose a \(\mathbb Z\)-basis of \(\Lambda\).  Its integral
endomorphisms form a discrete lattice in the real endomorphism
space, so a bounded set contains only finitely many of them.
Two nonnegative powers of \(g\) coincide.  Since \(g\) is
invertible, it has finite order on \(\Lambda\), and hence on
\(E^r\).  By PUN6 it acts as the identity modulo \(M\).

For completeness the required torsion argument uses only that
\(M\) is odd and at least three.  Choose an odd prime \(p\mid M\).
If a nonidentity integral matrix of finite order is the identity
modulo \(p\), a power has prime order \(q\) and is still of
the form \(1+p^a C\), where \(a\ge1\) is maximal and an entry
of \(C\) is not divisible by \(p\).  Expand its \(q\)-th power.
For \(q\ne p\), division by \(p^a\) and reduction modulo \(p\)
gives \(qC=0\).  For \(q=p\), division by \(p^{a+1}\) gives
\(C=0\pmod p\): every other term has an additional factor
\(p\), including the last because \(a(p-1)-1\ge1\).
Both are contradictions.  Thus \(g=1\), including when some
positive or negative signature multiplicity is zero.  In rank
zero the assertion is immediate.
\end{proof}
% END PUN8

% BEGIN PUN9
\paragraph{PUN9. One transvection modulus for all original integral isotropic vectors.}
\label{claim:PUN9}
For the same \(\Lambda,G,\xi_0\), with
\(\bar\xi_0=-\xi_0\), there is one integer \(m>0\) such that
for both \(\kappa_0=2\xi_0\) and \(\kappa_0=-2\xi_0\),
every \(v\in\Lambda\) with \(h_G(v,v)=0\), and every
\(s\in E\) with \(\bar s=s\) and \(s\) integral over
\(\mathbb Z\), the actual matrix
\(\operatorname{transvMat}(G,v,ms\kappa_0)\) belongs to the
original principal group.

\begin{proof}
Choose finite \(R\)-generators of \(\Lambda\).  Clear the
denominators of all their pairings multiplied by \(2\xi_0\),
using one positive rational integer \(d_0\).  Sesquilinearity
then gives
\(d_0(2\xi_0)h_G(\Lambda,\Lambda)\subseteq R\).
Set \(m=N d_0\).  For \(\alpha=ms\kappa_0\) define
\[
        T_\alpha(x)=x+\alpha h_G(x,v)v.
 \tag{PUN.10}
\]
This is exactly the original matrix: its \((a,b)\)-entry is
\(\delta_{ab}+\alpha v_a\sum_kG_{bk}\bar v_k\).
As \(\bar\alpha=-\alpha\) and \(h_G(v,v)=0\), expansion
gives \(h_G(T_\alpha x,T_\alpha y)=h_G(x,y)\), and
\(T_\alpha^{-1}=T_{-\alpha}\).  Because \(s\in R\), both
maps preserve \(\Lambda\), and
\((T_\alpha-1)\Lambda\subseteq N\Lambda\).
PUN6 puts them in the actual principal group.  The same \(m\)
works for both signs, every isotropic vector and every such \(s\);
it was chosen before these variables.
\end{proof}
% END PUN9

% BEGIN PUN10
\paragraph{PUN10. The whole original lattice-model supplier.}
\label{claim:PUN10}
Under the source's hypotheses for
\texttt{latticeModel\_of\_same\_signature\_hasseNorm}, every
datum \(D'\) with the same \((p,q)\) has an actual
\texttt{LatModel}\((l,D',\xi_0,G)\).  All its fields hold in
the prescribed space and on the original component lattice.

\begin{proof}
Apply PUN5 with \(\xi=-\xi_0\) to obtain the actual map
\(\phi':D'.L\hookrightarrow E^r\).  It has the original
perfect trace identity and full span.  Finite generation of
\(\Lambda'=\phi'D'.L\) permits a single nonzero integral
denominator for \(h_G(\Lambda',\Lambda')\), by clearing the
finitely many pairings of generators.  PUN7 supplies the exact
finite-index field; PUN8 supplies freeness for every original
flag frame and triangular multiplier; PUN9 supplies the one
positive transvection modulus with both original signs of
\(\kappa_0\).  These are precisely the fields of
\texttt{LatModel}; none is inferred merely from commensurability
of an auxiliary lattice.
\end{proof}
% END PUN10

% BEGIN PUN11
\paragraph{PUN11. Actual flag-domain models for any supplied family uniformization.}
\label{claim:PUN11}
Fix the original \(X_o,c,\sigma\), a \emph{supplied}
\(F_u:\texttt{FamilyUnif}\) on these objects, and the original
Hermitian coordinates \((G,\phi,\xi_0)\), with
\(\operatorname{Im}\tau\xi_0<0\) for every \(\tau\in\Psi\).
There is a \(\texttt{ModelsOf}\ F_u\ r\) satisfying the
whole original \texttt{ModelBase} and \texttt{NegXi}, and every
one of its matrices and imaginary scalars is exactly \(G\) and
\(\xi_0\).

\begin{proof}
The family record has finitely many component data \(D_i\),
each with \((D_i.p,D_i.q)=(D.p,D.q)\).  Apply PUN10 to each
component, choosing its injection \(\phi_i\) and put
\(\Lambda_i=\phi_iD_i.L\).  The inverse map onto its range
gives the required actual integral equivalence
\(D_i.L\simeq_R\Lambda_i\).

Define \(X_i\) using \(\Psi\) in (PUN.1), \(G\),
\(p_0=q\), \(\Lambda_i\), the principal group of PUN6, and
the constant character \(\chi(\lambda,g)=1\).  PUN3 applied
to \(-\xi_0\) gives the required signature.  PUN10 gives
finite generation, full span, the denominator, finite index and
freeness.  Take \(\kappa_0=2\xi_0\), whose imaginary part is
negative at every \(\tau\in\Psi\), and the modulus of PUN9.
All character laws hold identically, with finite-order exponent
one.  This constructs every field of the original \texttt{Data}.

The \texttt{ModelBase} trace and group fields are their defining
identities.  Its domain-stability field also holds literally.
Left multiplication by a unitary element preserves each Gram
matrix.  For right multiplication by an invertible upper
triangular matrix, the leading \(j\)-column Gram determinant
is multiplied by the positive factor
\(\prod_{a<j}|b_{aa}|^2\).  Thus every defining strict sign
inequality of \(\operatorname{GG}\) is preserved.  Rank,
CM type, and \texttt{NegXi} are the prescribed choices above.
This proves the exact arithmetic model producer; the existence
of \(F_u\) itself has not been assumed as a conclusion.
\end{proof}
% END PUN11

\subsection{The exact cusp-center lattice, in every cusp rank}

The next statements apply to \emph{every} original
\(X:\texttt{Data}\ E\ r\), not only to PUN11's models.
Fix an original cusp \(\mathrm{cu}:\texttt{CuspFD}\ X\ \rho\),
with linearly independent totally isotropic vectors
\(e_1,\ldots,e_\rho\), and write
\(Z(\beta)=\texttt{zMatE}\ X\ e\ \beta\).

% BEGIN PUN12
\paragraph{PUN12. The cusp-center parameter space and its faithful additive action.}
\label{claim:PUN12}
The original skew-Hermitian parameter space
\[
 \mathfrak z=\{\beta\in M_\rho(E):\bar\beta=-\beta^t\}
 \tag{PUN.11}
\]
has \(\mathbb Q\)-dimension \(\rho^2[E:\mathbb Q]/2\).
The original matrices satisfy
\[
 Z(0)=1,\quad Z(\beta+\beta')=Z(\beta)Z(\beta'),\quad
 Z(\beta)^{-1}=Z(-\beta),
 \tag{PUN.12}
\]
and \(\beta\mapsto Z(\beta)\) is injective.
For \(\beta\in\mathfrak z\), \(Z(\beta)\) is an isometry
of the original Hermitian form.

\begin{proof}
Choose a nonzero anti-invariant \(w\in E\).  Then
\(E=F\oplus Fw\).  Each diagonal entry of \(\beta\) lies in
\(Fw\), and each pair above and below the diagonal is determined
by one arbitrary element of \(E\).  This gives dimension
\(\rho[F:\mathbb Q]+\rho(\rho-1)[E:\mathbb Q]/2
=\rho^2[E:\mathbb Q]/2\).

In the source's convention the matrix acts by
\[
 Z(\beta)x=x+\sum_a\Big(\sum_b\beta_{ab}h_G(x,e_b)\Big)e_a.
 \tag{PUN.13}
\]
Every displacement is in the span of the \(e_a\), and all
functionals \(h_G(-,e_b)\) vanish on that span.  Consequently
any product of two displacement operators is zero, proving
(PUN.12).  Nondegeneracy of \(h_G\) and independence of the
\(e_b\) make \(x\mapsto(h_G(x,e_b))_b\) surjective.
If the displacement is zero, independence of the \(e_a\)
therefore gives \(\beta=0\).  Finally, expand the pairing of
the two sides of (PUN.13).  The term involving two
displacements vanishes by total isotropy.  The two remaining
cross terms cancel precisely because
\(\beta_{ab}+\bar\beta_{ba}=0\).  This proves the isometry
claim with the actual matrix and its actual order.
\end{proof}
% END PUN12

% BEGIN PUN13
\paragraph{PUN13. Finite generation and full rational span of the actual center subgroup.}
\label{claim:PUN13}
The original subgroup
\[
 Z_\Gamma=\texttt{zGam}\ X\ \mathrm{cu}
       =\{\beta\in\mathfrak z:Z(\beta)\in\Gamma\}
 \tag{PUN.14}
\]
is a finitely generated abelian group, and every
\(\beta\in\mathfrak z\) has a positive integral multiple in
\(Z_\Gamma\).  In particular its rational span is exactly
\(\mathfrak z\).

\begin{proof}
PUN12 shows that (PUN.14) is an additive subgroup.  Choose finite
\(R\)-generators \(v_1,\ldots,v_t\) of the original
\(\Lambda\), and map
\[
 \beta\longmapsto
        ((Z(\beta)-1)v_1,\ldots,(Z(\beta)-1)v_t).
 \tag{PUN.15}
\]
This is additive and injective: the \(v_i\) span \(E^r\),
and PUN12 is faithful.  For \(\beta\in Z_\Gamma\) its
image lies in \(\Lambda^t\).  The latter is finite free over
\(\mathbb Z\); all its subgroups are finitely generated.
This proves finite generation without first assuming a lattice
or a bounded fundamental domain for the center.

For any \(\beta\in\mathfrak z\), clear the finitely many
denominators of \((Z(\beta)-1)v_i\) relative to \(\Lambda\).
There is a positive integer \(d_0\) with
\(d_0(Z(\beta)-1)\Lambda\subseteq\Lambda\).
By (PUN.12), \(Z(jd_0\beta)\) and its inverse preserve
\(\Lambda\) for every integer \(j\), and PUN12 says they
are unitary.  The original \texttt{Gamma\_index} field gives
finitely many cosets of \(\Gamma\) in the unitary lattice
stabilizer.  Two of the nonnegative powers
\(Z(jd_0\beta)\) have the same coset.  Their quotient is
\(Z((j-i)d_0\beta)\in\Gamma\), for some \(j>i\).
Thus \((j-i)d_0\beta\in Z_\Gamma\), proving the required
positive multiple and the full rational span.  Neither normality
of \(\Gamma\) nor a uniform congruence level was used.
\end{proof}
% END PUN13

% BEGIN PUN14
\paragraph{PUN14. The literal all-rank \texttt{LatZ} producer.}
\label{claim:PUN14}
For every CM number field \(E\), every \(r,\rho\in\mathbb N\),
every component type \(I\), and every actual family
\(X:I\to\texttt{Data}\ E\ r\), the original
\(\texttt{FJ.LatZ}\ X\ \rho\) holds.

\begin{proof}
For each original cusp of each component, PUN13 identifies the
actual \(Z_\Gamma\) as a finitely generated subgroup of a
\(\mathbb Q\)-vector space.  It is torsion free, so the
structure theorem for finitely generated abelian groups gives
a \(\mathbb Z\)-basis.  Its length is the dimension of its
rational span, hence by PUN12--PUN13 it is exactly
\(\rho^2[E:\mathbb Q]/2\).  The original CM-type condition
pairs the embeddings into two complementary equally large sets,
so \(|X.\Psi|=[E:\mathbb Q]/2\).  The basis length is therefore
the literal \(\texttt{nZ}\ X\ \rho=|X.\Psi|\rho^2\).
Regard the basis vectors as the
original skew-Hermitian matrices.  For every matrix \(\beta\),
membership in \(\texttt{zGam}\) is equivalent to membership
in their \(\mathbb Z\)-span, and they are \(\mathbb Z\)-linearly
independent.  These are precisely the two fields of the original
\texttt{IsZBasis}.  The construction can be chosen for each cusp,
which is exactly \texttt{LatZ}'s quantifier order.  If \(\rho=0\)
the zero group has its empty basis; if no such cusp exists the
corresponding assertion is vacuous.  No rank or component
nonemptiness restriction has been introduced.
\end{proof}
% END PUN14

\subsection{Periods and the level-one exponent dictionary}

% BEGIN PUN15
\paragraph{PUN15. Periods for every original ordered type function.}
\label{claim:PUN15}
For every original datum \(D\), and every original
\(\mathrm{ty}:\{1,\ldots,g\}\to\operatorname{Hom}(E,\mathbb C)\)
with \(\texttt{IsTypeFn}(H,p,q,\mathrm{ty})\),
\[
 g=[E:\mathbb Q](p+q)/2,
 \qquad \exists P:L\longrightarrow\mathbb C^g,
                           \texttt{IsPeriod}(D,\mathrm{ty},P).
 \tag{PUN.16}
\]
Such type functions exist.  This proves the two unconditional
period-existence clauses of the original
\texttt{typeFn\_period\_exist}, without assuming \texttt{N1c}.

\begin{proof}
Let \(V=L\otimes_{\mathbb Z}\mathbb R\), and give it the
complex structure supplied in \texttt{Datum.dom}, declaring
\(i v=Jv\).  The \(E\)-action is complex linear.  As \(E\)
is separable, its complexification splits into its embedding
eigenspaces.  For \(a\in E\) the original trace identities
give
\[
 \operatorname{tr}_{\mathbb C}(a\mid V,J)
 =\tfrac12\bigl(\operatorname{tr}_{\mathbb R}a
                  -i\operatorname{tr}_{\mathbb R}(Ja)\bigr)
 =\sum_{\tau\in\Psi}
                       \bigl(p\tau(a)+q\bar\tau(a)\bigr).
 \tag{PUN.17}
\]
The distinct embeddings are linearly independent, so the
\(\tau\)-eigenspace has exactly the prescribed dimension.
Choosing a basis in each eigenspace, and ordering it according
to the supplied \(\mathrm{ty}\), gives an \(E\)-equivariant
complex-linear isomorphism \(T:(V,J)\simeq\mathbb C^g\).
It also gives the dimension formula; choosing any ordering
proves the existence of a type function.

Define \(P(x)=T(1\otimes x)\).  This is additive, injective,
and has the exact coordinate action
\(P(ax)_j=\mathrm{ty}_j(a)P(x)_j\).  A \(\mathbb Z\)-basis
of \(L\) is a real basis of \(V\), so \(P(L)\) is discrete
and spans \(\mathbb C^g\) over \(\mathbb R\).
On \(V\) define
\[
                 H(u,v)=B(Ju,v)+iB(u,v).
 \tag{PUN.18}
\]
The relations \(J^2=-1\), \(B(Ju,Jv)=B(u,v)\) and
alternation show that \(H\) is Hermitian and complex linear
in the first variable.  Its real diagonal value is
\(B(Ju,u)>0\) for \(u\ne0\).  Transporting it through \(T\)
gives a positive Hermitian matrix \(H_m\) satisfying
\(\operatorname{Im}H_m(Px,Py)=B(x,y)\).  These are all the
literal fields of \texttt{IsPeriod}.  The construction also
works for the empty type and zero lattice.
\end{proof}
% END PUN15

% BEGIN PUN16
\paragraph{PUN16. An adapted cusp pseudobasis and surjectivity of the exact pairing map.}
\label{claim:PUN16}
Assume the original \(\Lambda\subset E^r\) has the perfect
pairing \(B(x,y)=\operatorname{Tr}_{E/\mathbb Q}(-\xi_0h_G(x,y))\).
For any rational totally isotropic subspace \(W\), put
\(C=\Lambda\cap W\).  It has a pseudobasis
\[
                    C=\bigoplus_{j=1}^{\rho}\mathfrak a_j e_j,
                    \qquad \mathfrak a_j\ne0,
 \tag{PUN.19}
\]
with \(e_j\) an \(E\)-basis of \(W\).  For this basis the
map
\[
 \Lambda\longrightarrow\prod_j
        \mathfrak D^{-1}\bar{\mathfrak a}_j^{-1},
 \quad x\longmapsto(-\xi_0h_G(x,e_j))_j,
 \qquad\mathfrak D=\mathfrak D_{E/\mathbb Q},
 \tag{PUN.20}
\]
is onto.

\begin{proof}
The submodule \(C\) is finite and torsion free over the Dedekind
domain \(R\), hence projective.  It spans \(W\): an arbitrary
vector of \(W\) has a nonzero integral multiple in \(\Lambda\)
by fullness of \(\Lambda\).  To obtain (PUN.19), choose a
nonzero coordinate functional on \(W\).  Its image on \(C\)
is a nonzero fractional ideal, hence projective, so the exact
sequence onto that ideal splits.  The kernel has rank one less.
Inducting gives the direct sum of fractional ideals and its
corresponding rational basis.  This argument permits nonprincipal
ideals and does not make \(C\) free over \(R\).

Moreover \(C\) is a primitive \(\mathbb Z\)-sublattice of
\(\Lambda\): if \(n x\in C\) with \(n\ne0\) and
\(x\in\Lambda\), then \(x\in W\).  Thus
\(\Lambda/C\) is finite free over \(\mathbb Z\), and every
integral functional on \(C\) extends to \(\Lambda\).
For any tuple \(t_j\in\mathfrak D^{-1}\bar{\mathfrak a}_j^{-1}\),
the formula
\[
          \sum_j a_j e_j\longmapsto
                  \operatorname{Tr}_{E/\mathbb Q}
                                      \bigl(\sum_j t_j\bar a_j\bigr)
 \tag{PUN.21}
\]
is integral on \(C\).  Extend it to \(\Lambda\) and use
perfection to represent the extension by \(B(x,-)\).
Testing (PUN.21) with one coefficient in \(\mathfrak a_j\)
at a time and using nondegeneracy of the trace pairing shows
\(-\xi_0h_G(x,e_j)=t_j\) for all \(j\).
Conversely these coordinates necessarily lie in the displayed
ideals by integrality on \(C\).  This proves the exact image,
not only its rational span.
\end{proof}
% END PUN16

% BEGIN PUN17
\paragraph{PUN17. The actual center congruence is an entrywise fractional-ideal condition.}
\label{claim:PUN17}
Suppose \(\Gamma\) is the original principal group and \(A=1\).
In an adapted cusp basis as in (PUN.19), set
\[
            I_{ij}=M\mathfrak a_i\mathfrak D\bar{\mathfrak a}_j.
 \tag{PUN.22}
\]
For a skew-Hermitian matrix \(\beta\),
\[
 \beta\in Z_\Gamma
       \quad\Longleftrightarrow\quad
           \beta_{ij}/\xi_0\in I_{ij}\quad\hbox{for every }i,j.
 \tag{PUN.23}
\]

\begin{proof}
PUN6 identifies \(\Gamma\) with the full congruence kernel
modulo \(M\).  A displacement \((Z(\beta)-1)x\) lies in
\(W\).  Since \(M\Lambda\cap W=MC\), the congruence
condition is exactly
\[
  \sum_j\beta_{ij}h_G(x,e_j)\in M\mathfrak a_i
                      \quad(x\in\Lambda,\ 1\le i\le\rho).
 \tag{PUN.24}
\]
PUN16 says that \(-\xi_0h_G(x,e_j)\) ranges independently
over \(\mathfrak D^{-1}\bar{\mathfrak a}_j^{-1}\).
Consequently (PUN.24) is equivalent, entry by entry, to
\[
 (\beta_{ij}/\xi_0)
        \mathfrak D^{-1}\bar{\mathfrak a}_j^{-1}
                                      \subseteq M\mathfrak a_i,
\]
which is exactly (PUN.23), by invertibility of nonzero fractional
ideals.  Conversely the same calculation proves the congruence
for every \(x\); the inverse has displacement with the opposite
sign and hence preserves \(\Lambda\) too.  PUN12 supplies
unitarity.  Thus no stabilizer condition was lost in reducing
to entries.  Notice \(I_{ji}=\bar I_{ij}\), because the
different is conjugation stable.
\end{proof}
% END PUN17

% BEGIN PUN18
\paragraph{PUN18. The factor of two in the exact Fourier phase.}
\label{claim:PUN18}
Let \(\chi\) be Hermitian and let \(I_{ij}\) be as in
(PUN.22).  The phase in the original \texttt{prE} satisfies
\[
 \bigl(\operatorname{prE}(\chi,\beta)\in\mathbb Z
                 \text{ for every }\beta\in Z_\Gamma\bigr)
 \quad\Longleftrightarrow\quad
 \operatorname{Tr}_{E/\mathbb Q}(\chi_{ij}s)\in\mathbb Z
                 \text{ for every }s\in I_{ij},\ i,j.
 \tag{PUN.25}
\]
This includes the original choices \(\kappa_0=\pm2\xi_0\).

\begin{proof}
Set \(S=\beta/\xi_0\).  Then \(S\) is Hermitian, and PUN17
is the condition \(S_{ij}\in I_{ij}\).  The definition of
\texttt{prE}, summing the real parts over one CM type, gives
\[
 \operatorname{prE}(\chi,\beta)
       =\pm\frac12\operatorname{Tr}_{E/\mathbb Q}
                              \Bigl(\sum_{i,j}\chi_{ij}S_{ij}\Bigr).
 \tag{PUN.26}
\]
The sign is immaterial for integral membership.  If \(i\ne j\),
test (PUN.26) on the matrix with entries \(s,\bar s\) in
positions \((i,j),(j,i)\).  Its half-trace is
\(\operatorname{Tr}_{E/\mathbb Q}(\chi_{ij}s)\).
For a diagonal position test the real entry \(s+\bar s\).
Because \(\chi_{ii}\in F\), its half-trace is the same
quantity.  This proves the forward implication even at primes
above \(2\).

Conversely the right side implies
\(\chi_{ij}I_{ij}\subseteq\mathfrak D^{-1}\), since
each ideal is closed under multiplication by \(R\).
For an admissible \(S\) the sum
\(u=\sum_{i,j}\chi_{ij}S_{ij}\) lies in
\(\mathfrak D^{-1}\), and conjugating and interchanging the
indices gives \(\bar u=u\).  The original host element
\(w_0\in R\) has \(w_0+\bar w_0=1\), by
\(2w_0=1+\mathrm{rt}\).  Therefore
\[
 \tfrac12\operatorname{Tr}_{E/\mathbb Q}(u)
       =\operatorname{Tr}_{E/\mathbb Q}(w_0u)\in\mathbb Z.
 \tag{PUN.27}
\]
This proves integrality of (PUN.26).  The trace-one element,
not an unjustified division by \(2\) in the integer ring,
is what supplies the required even trace.
\end{proof}
% END PUN18

% BEGIN PUN19
\paragraph{PUN19. The whole original exponent match when \(l.A=1\).}
\label{claim:PUN19}
For every supplied family \(F_u\), every \texttt{ModelsOf}
satisfying the original \texttt{ModelBase}, and every
\(\rho\in\mathbb N\),
\[
                  l.A=1\quad\Longrightarrow\quad
                         \texttt{FJ.ExpoMatch}\ M.X\ 1\ \rho\ l.M.
 \tag{PUN.28}
\]
Thus the asserted Fourier exponents use the original coefficient
lattice \(\texttt{XcFD}\ 1\ \mathrm{cu}\), every original
cusp basis, and the exact modulus \(l.M\).

\begin{proof}
In the adapted basis of PUN16, the exact definitions of
\(\texttt{crdI}\), \(\texttt{dualZE}\) and
\(\texttt{XcFD}\) give
\[
 \begin{split}
 \texttt{crdI}(1,\mathrm{cu})&=\prod_i\mathfrak a_i,\\
 \texttt{XcFD}(1,\mathrm{cu})&=
                      \prod_i\mathfrak D^{-1}\mathfrak a_i^{-1},\\
 \texttt{dualOE}(\texttt{XcFD}(1,\mathrm{cu}))&=
                      \prod_i\mathfrak D\mathfrak a_i.
 \end{split}                                                   \tag{PUN.29}
\]
The first equality follows by expanding a vector in the
pseudobasis; the second is the trace-dual calculation of PCL2;
the third is ordinary fractional-ideal duality.  The trace dual
is already an \(R\)-module, so the additional span in
\texttt{XcFD} does not enlarge it.

For a nonzero fractional ideal \(J\),
\(\operatorname{Tr}(zJ)\subseteq\mathbb Z\) is equivalent
to \(zJ\subseteq\mathfrak D^{-1}\).  With \(J=I_{ij}\),
the right side of (PUN.25) is therefore equivalent to
\[
     M\chi_{ij}(\mathfrak D\mathfrak a_i)
                     (\mathfrak D\bar{\mathfrak a}_j)\subseteq R
                            \qquad\text{for all }i,j.
 \tag{PUN.30}
\]
Testing vectors supported in one coordinate and then adding
shows that this is precisely
\[
    M\sum_{i,j}u_i\chi_{ij}\bar v_j\in R
       \quad\text{for all }u,v\in
                     \texttt{dualOE}(\texttt{XcFD}(1,\mathrm{cu})).
 \tag{PUN.31}
\]
PUN18 proves equivalence with integral phase on the actual
\(\texttt{zGam}\), which is the required statement in this
basis.

To verify the original arbitrary-basis quantifier, write the
given cusp basis as \(e_a^0=\sum_b P_{ab}e_b\), with
\(P\in\operatorname{GL}_\rho(E)\).  Direct substitution in
(PUN.13) gives
\[
 Z_{e^0}(\beta)=Z_e(P^t\beta\bar P).
 \tag{PUN.32}
\]
The corresponding changes of the Hermitian argument and dual
coordinates are \(\chi^0=P\chi\bar P^t\) and
\(u=u^0P\), \(v=v^0P\).  The identities
\[
 \operatorname{tr}((\chi^0)^t\beta)
       =\operatorname{tr}(\chi^tP^t\beta\bar P),
 \qquad
 u^0\chi^0\bar v^{0t}=(u^0P)\chi\overline{(v^0P)}^{\,t}
 \tag{PUN.33}
\]
preserve the phase and the integral expression.  The coordinate
pairing shows that \(\texttt{dualOE}(\texttt{XcFD})\) transforms
by exactly this row-vector map; using \(P^{-1}\) gives both
directions.  Hence (PUN.31) and the phase condition transport
back to the given cusp, not to a substitute basis.  The empty
rank case has zero sums and is included.  Applying this to
each component proves the literal all-\(\rho\) endpoint.
\end{proof}
% END PUN19

\subsection{Actual moduli points from integral periods}

\paragraph{Classical complex-geometric foundations used in this subsection.}
We use the Riemann algebraicity criterion for a complex torus,
Appell--Humbert classification and its positivity criterion,
GAGA for proper complex schemes, and the representing dual abelian variety with its
Poincare bundle.  The universal property of the dual is used for
\emph{all} test schemes.  These are named external foundations,
not numbered reconstructions of their classical proofs.
The analytic references are Conrad, \emph{Abelian Varieties:
Geometry, Parameter Spaces, and Arithmetic}, Theorems 2.2.3,
3.1.3 and 4.3.2, and Example 4.1.4; the all-scheme universal
property is recorded in Milne, \emph{Abelian Varieties},
\S8, especially Remark 8.7 and the paragraph preceding
Remark 8.15.\footnote{Author texts:
\href{https://math.stanford.edu/~conrad/papers/cmchina.pdf}{Conrad's notes}
and \href{https://www.jmilne.org/math/CourseNotes/AV.pdf}{Milne's notes}.
The metadata lists each foundation separately.  None of these
references is being cited for the original \texttt{N1c} record.}
The explicit CM line, commutator sign, integral markings and
literal moduli installation are proved next.

% BEGIN PUN20
\paragraph{PUN20. The canonical primitive CM line on a period torus.}
\label{claim:PUN20}
Let \(P:L\to\mathbb C^g\) be an original \texttt{IsPeriod}
for \(D\), and let \(H_m\) be its positive Hermitian form.
The quotient \(T=\mathbb C^g/P(L)\) has an algebraic abelian
variety model with the given \(R\)-action and a rigidified,
symmetric, cubical, ample line bundle \(\mathcal L\) such that
the exact \texttt{SAB.adj}, \texttt{SAB.prim} and
\texttt{Out.KTriv} fields hold.  The exponential map is the
original quotient map, with its actual kernel \(P(L)\), and
its tangent map is an isomorphism in both directions in the
source's dual-number sense.

\begin{proof}
For \(\lambda=P(a)\) put
\[
       \alpha(\lambda)=\exp\bigl(\pi i B(a,w_0a)\bigr)\in\{1,-1\}.
 \tag{PUN.34}
\]
The action of \(w_0\) is well defined on the lattice.  Using
\(w_0+\bar w_0=1\) and the adjoint identity shows
\[
 B(a+b,w_0(a+b))-B(a,w_0a)-B(b,w_0b)
                                  \equiv B(a,b)\pmod2.
\]
Thus \(\alpha\) is the semicharacter for the alternating
form \(\operatorname{Im}H_m\).  The factors
\[
 j(\lambda,z)=\alpha(\lambda)
       \exp\bigl(\pi H_m(z,\lambda)+\pi H_m(\lambda,\lambda)/2\bigr)
 \tag{PUN.35}
\]
satisfy the lattice cocycle law, as direct multiplication
shows.  They define an analytic line bundle on \(T\).
The named Riemann and Appell--Humbert theorems make \(T\)
projective and this line ample; GAGA supplies the algebraic
variety, bundle, group law and all the \(R\)-endomorphisms.
The quotient's canonical trivial frame at \(0\) fixes its
rigidification.  Since \(\alpha(-\lambda)=\alpha(\lambda)\),
the bundle is symmetric.

Here are the original additional line identities.  Write
\(\mathcal B(\mathcal M)=m^*\mathcal M\otimes
p_1^*\mathcal M^{-1}\otimes p_2^*\mathcal M^{-1}\otimes
e^*\mathcal M\) for the normalized difference bundle.
For an Appell--Humbert bundle with form \(H_m\), the pullback
\((1,w_0)^*\mathcal B(\mathcal M)\) has form
\[
 H_m(u,w_0v)+H_m(w_0u,v)=H_m(u,v)
\]
and semicharacter exactly (PUN.34), independently of the
semicharacter of \(\mathcal M\).  Indeed the semicharacter
ratio is
\(\alpha_{\mathcal M}(a+w_0a)\alpha_{\mathcal M}(a)^{-1}
\alpha_{\mathcal M}(w_0a)^{-1}=\exp(\pi iB(a,w_0a))\).
Taking \(\mathcal M=\mathcal L\) proves the literal primitive
identity.  The adjoint relation for \(B\), and hence for
\(H_m\), proves
\((a,1)^*\mathcal B(\mathcal L)
\simeq(1,\bar a)^*\mathcal B(\mathcal L)\) for every
\(a\in R\): their mixed Hermitian forms and their mixed
semicharacters agree.  This is precisely \texttt{SAB.adj}.
The third finite difference of these quadratic factors cancels.
It gives the normalized cubical trivialization; its coherence
identities hold because the ratios are global holomorphic units
on powers of the compact connected torus and are one along the
normalizing axes.  GAGA transfers these identities, not only
their Chern classes, to algebraic line-bundle isomorphisms.

Perfection of \(B\) identifies its integral lattice with its
dual.  The polarization map \(T\to T^\vee\) associated to
\(\mathcal L\) is therefore an analytic isomorphism; GAGA
makes it an algebraic isomorphism.  For an arbitrary commutative
\(\mathbb C\)-algebra \(A_0\), a point \(x\) whose
translation-difference bundle is locally trivial over
\(\operatorname{Spec}A_0\) maps to zero in the representing
dual abelian variety.  The isomorphism just proved forces
\(x=0\).  This is the full \texttt{KTriv} test, including
nonreduced algebras, rather than only a statement about complex
points.  Algebraic ampleness supplies \texttt{RelAmple} in its
printed form: choose a very ample power, use its projective
coordinate sections, and on each nonzero locus their ratios
generate the affine coordinate ring.

Finally the quotient map is holomorphic and locally biholomorphic.
Algebraic regular functions are analytic on their open domains,
giving the exact \texttt{HolUnif.hol} condition.  Analytic
coordinates at the origin identify its derivative with an
isomorphism of complex tangent spaces.  A tangent vector of a
smooth complex scheme is the derivation, or equivalently the
map from the dual numbers, evaluating every local function at
its directional derivative.  This is exactly
\texttt{IsDiff}/\texttt{IsDerivAt0}; both required uniqueness
clauses follow.  Reading the pullback frame in any algebraic
bundle atlas gives the exact \texttt{HolUnif.fr} and
\texttt{IsAH} identities (PUN.35), normalized at the origin.
\end{proof}
% END PUN20

% BEGIN PUN21
\paragraph{PUN21. The original negative Weil-pairing sign.}
\label{claim:PUN21}
For the algebraic object of PUN20, every \(n>0\) and
\(a,b\in L\), the original \texttt{Out.WeilIs} relation on
\(e(Pa/n),e(Pb/n)\) has scalar
\[
                   \exp\bigl(-2\pi i B(a,b)/n\bigr).
 \tag{PUN.36}
\]
It remains this scalar after every field embedding
\(\mathbb C\to k\).

\begin{proof}
For \(x=Pa/n\), a generator of the source's
\(\mathcal L_z^{\otimes n}\otimes
\mathcal L_{z+x}^{\otimes(-n)}\) is represented on the
covering space by
\[
                         a_x(z)=\exp(-\pi n H_m(z,x)).
 \tag{PUN.37}
\]
Its lattice transformation is the correct quotient of the
factors (PUN.35): the apparent difference between
\(H_m(\lambda,x)\) and \(H_m(x,\lambda)\) has exponential
\(\exp(2\pi iB(\lambda,Pa))=1\).
It descends to a nowhere-zero analytic section, which is
algebraic by GAGA.  In the exact multiplication order of
\texttt{CommIsOf}, the commutator is
\[
 \frac{a_x(z)a_y(z+x)}{a_y(z)a_x(z+y)}
      =\exp\bigl(-\pi n(H_m(x,y)-H_m(y,x))\bigr)
      =\exp\bigl(-2\pi i B(a,b)/n\bigr).
 \tag{PUN.38}
\]
Every other pair of generators differs by nonzero constants,
because the abelian variety is proper and connected; these
constants cancel.  Thus the definition's universal condition on
all generators holds, not merely a computation for a chosen
pair.  The exhibited isomorphisms and their scalar equality
are algebraic.  Pulling them back along \(\mathbb C\to k\)
preserves the identities and the scalar is its image in \(k\).
Proper geometrically connected base change gives the same
constant-ratio assertion over \(k\).  This verifies the literal
local-trivialization definition of \texttt{WeilIs}.
\end{proof}
% END PUN21

% BEGIN PUN22
\paragraph{PUN22. Every original marked period gives an actual moduli point.}
\label{claim:PUN22}
The whole original \texttt{exists\_point\_of\_isPeriod} is
supplied by PUN20--PUN21 and the original \texttt{Moduli.rep}.
In particular retain all its data: \(D',t,P\), the compatible
maps \(f_n:D.L\to D'.L\), multipliers \(\nu_n\), the fixed
\(\sigma:K_A(l)\to\mathbb C\), the hypotheses
\(D'.p=D.p\), \(D'.q=D.q\), \(p+q\ge1\),
\(\sigma(\rho)=\mathrm{rt}_0\), the given ordered
\(\texttt{IsTypeFn}\), the given \(\texttt{IsPeriod}\), and
\[
 \begin{gathered}
 f_{nm}\equiv f_n\pmod n,\qquad
 f_n:D.L/nD.L\simeq_R D'.L/nD'.L,\\
 B'(f_na,f_nb)\equiv\nu_nB(a,b)\pmod n,\qquad
 \sigma(\zeta_A)=\exp(-2\pi i\nu_A/A),\qquad t=f_N.
 \end{gathered}                                                   \tag{PUN.39}
\]
There is a point \(y\) of the supplied \(X_o\) over
\(\sigma\), with the source's actual exponential map,
kernel, tangent bijections, \(R\)-action, all Weil pairings,
and both printed level identities.

\begin{proof}
Construct the abelian variety and its primitive line from
\(D',P\) by PUN20.  For \(n>0\) define
\[
 v_n(a)=e(P(f_na)/n),\qquad
                         \zeta_n=\exp(-2\pi i\nu_n/n).
 \tag{PUN.40}
\]
The induced map on \(L/nL\) is an isomorphism, and the
quotient description identifies all \(n\)-torsion points with
\(\tfrac1nP(L')/P(L')\).  Hence \(v_n\) has the exact
kernel \(nL\) and the exact torsion image.  Addition and
\(R\)-linearity are immediate.  The first congruence in
(PUN.39) gives the tower relation.

The multiplier \(\nu_n\) is a unit modulo \(n\): the two
perfect alternating forms modulo \(n\), transported by the
isomorphism \(f_n\), differ by this multiplier.  Since the
rank is positive, an integral symplectic pair with pairing one
tests its invertibility.  The same pair and the congruences
in (PUN.39) give \(\nu_{nm}\equiv\nu_n\pmod n\).
Thus \(\zeta_n\) is primitive and
\(\zeta_{nm}^m=\zeta_n\).  PUN21 gives
\(e_n(v_n(a),v_n(b))=\zeta_n^{B(a,b)}\).
For \(n=0\), set \(v_0=1\), \(\zeta_0=1\); no positive-level
clause is imposed there.  These formulas construct every
field of the original \texttt{TateIso}.

The same construction is valid at every algebraically closed
field \(k\) and every \(\beta:\mathbb C\to k\), with no
continuity condition on \(\beta\).  Indeed \([n]\) is an
algebraic finite etale isogeny in characteristic zero: on the
complex torus it is a finite covering of degree \(n^{2g}\),
and GAGA gives finiteness and the invertible differential.
Its kernel is therefore the constant finite etale scheme on
the complex torsion points and remains so after base change.
All torsion points over \(k\) are exactly their scalar
extensions.  The eigenbasis of the complex Lie algebra also
base changes; every embedding of the number field into \(k\)
is an extension through \(\beta\) of one of its complex
embeddings.  It therefore gives the precise \texttt{DetCond}
for every \(k\)-valued point in \texttt{Obj.geom}.

Choose \(i\in\mathfrak b\), \(j\in\bar{\mathfrak b}\)
with \(i+j=1\).  Define the half-level marking by
\(\eta(a)=v_M(ja)\), and define \(\eta_A(a)=v_A(a)\).
The first kills \(\mathfrak bL\), since
\(j\mathfrak b\subseteq MR\), and it equals \(v_M(a)\)
on \(\bar{\mathfrak b}L\), since
\(i\bar{\mathfrak b}\subseteq MR\).  The second kills
\(AL\).  Along with PUN20, these give every field of the
original \texttt{Obj}, including its full geometric test
over every \(\beta\).  Its base values are
\(\sigma(\rho)\) and \(\sigma(\zeta_A)\).

Apply the given representing property to this object over
the reduced scheme \(\operatorname{Spec}\mathbb C\).
It gives the actual \(y:\operatorname{Spec}\mathbb C\to X_o\)
and an isomorphism with the pullback of the universal object.
Transport the quotient exponential map through this
isomorphism.  Its group, holomorphic, tangent, kernel and
Weil fields are unchanged.  Finally, for \(n=M,A\), the
congruence \(f_N\equiv f_n\pmod n\) gives
\(e(P(t a)/n)=v_n(a)\).  On the prescribed half-lattice
this is exactly \(\eta(a)\), and for all \(a\) it is
\(\eta_A(a)\) at level \(A\).  This establishes both
literal level identities at the specified point.
\end{proof}
% END PUN22

% BEGIN PUN23
\paragraph{PUN23. Equality of actual period points is exactly marked integral isometry.}
\label{claim:PUN23}
Take two sets of original period-point data
\((D',t,P,y,e)\) and \((D'',t',P',y',e')\) satisfying the
full conclusions used in \texttt{point\_eq\_iff\_isometry\_closed}.
Then \(y=y'\) if and only if there are an integral
\(R\)-linear equivalence \(\gamma:D'.L\simeq D''.L\)
and a complex-linear equivalence \(h:\mathbb C^g\simeq\mathbb C^g\)
such that
\[
 \begin{gathered}
 B''(\gamma a,\gamma b)=B'(a,b),\qquad P'\gamma=hP,\\
 \gamma t(a)-t'(a)\in MD''.L\quad(a\in\bar{\mathfrak b}D.L),\\
 \gamma t(a)-t'(a)\in AD''.L\quad(a\in D.L).
 \end{gathered}                                                   \tag{PUN.41}
\]

\begin{proof}
If \(y=y'\), the two exponential maps are holomorphic covering
homomorphisms of the same complex torus, with invertible
derivatives.  Lifting from the simply connected space
\(\mathbb C^g\), or using the universal cover of a complex
Lie group, gives a unique additive holomorphic equivalence
\(h\) with \(e'h=e\).  An additive holomorphic map is
complex linear.  The kernel identities give
\(hP(L')=P'(L'')\), hence the integral equivalence
\(\gamma\).  Compatibility with the \(R\)-action implies
its \(R\)-linearity.  Comparing the Weil scalars for every
positive \(n\) shows
\(B''(\gamma a,\gamma b)-B'(a,b)\in n\mathbb Z\) for every
\(n\); it is therefore zero.  Comparing the two specified
level markings gives exactly the last two congruences of
(PUN.41).

Conversely (PUN.41) gives an analytic isomorphism of complex
tori, which is algebraic by GAGA and respects the \(R\)-action.
The given Weil-pairing identities for all \(n\) identify each
fiber line's integral Riemann form with \(B'\) or \(B''\),
by the same divisibility argument.  The primitive condition
on an original \texttt{SAB} then determines its semicharacter:
the calculation in PUN20 says that
\((1,w_0)^*\mathcal B(\mathcal M)\) has precisely (PUN.34).
Thus the isomorphism preserves the actual primitive line,
not just its polarization class.  Normalize the line
isomorphism at zero; the cubical isomorphism follows by
uniqueness of the normalized cube trivialization.  The two
congruences in (PUN.41) preserve \(\eta\) and \(\eta_A\).
We have an isomorphism of the original moduli objects.
Uniqueness in the supplied \texttt{Moduli.rep}, applied to
\(\operatorname{Spec}\mathbb C\), forces \(y=y'\).
This proves both directions of the literal endpoint without
restricting to a chosen representative family.
\end{proof}
% END PUN23

% BEGIN PUN24
\paragraph{PUN24. A marked genus exists for every prescribed base embedding.}
\label{claim:PUN24}
For every original \(\sigma:K_A(l)\to\mathbb C\) satisfying
\(\sigma(\rho)=\mathrm{rt}_0\), the original
\texttt{InGenus}\((D,l,\sigma,D',t)\) is inhabited.
One can take \(D'=D\).  Together with PUN15 and PUN22, this
constructs a point of the original moduli space over that
\(\sigma\) whenever \(p+q\ge1\).

\begin{proof}
Write \(\sigma(\zeta_A)=\exp(-2\pi i k/A)\), with
\(k\) a unit modulo \(A\).  Choose
\(u\in\widehat{\mathbb Z}^{\times}\) reducing to \(k\)
modulo \(A\): at primes dividing \(A\) use the indicated
unit residue, and choose one at the other primes.
At every finite place of \(F\), the original unramified
quadratic integer algebra has surjective norm on units by
MF2, including at \(2\).  Therefore choose
\(s\in(\mathcal O_E\otimes\widehat{\mathbb Z})^\times\)
with \(s\bar s=u\).  The choices are independent at the
finite places and combine in this profinite product.

For each \(n>0\), choose \(s_n\in R\) lifting the residue
of \(s\) in \(R/nR\), and choose an integer \(\nu_n\)
lifting \(u\bmod n\).  Put \(f_n(a)=s_na\).
Reduction is invertible, \(f_{nm}\equiv f_n\pmod n\), and
the adjoint identity gives
\[
 B(f_na,f_nb)=B(a,\bar s_ns_nb)
                              \equiv\nu_nB(a,b)\pmod n.
\]
Take \(t=f_N\).  The required root-of-unity identity is the
chosen one at \(A\).  These are all fields of the original
\texttt{InGenus}; the omitted \(n=0\) values may be chosen
arbitrarily.  If \(A=1\) the congruence imposes no restriction
and choose \(u=1\).  The construction neither changes the
coefficient embedding nor assumes a global integral norm
solution.  PUN15 supplies a period of \(D\) for each desired
type, and PUN22 gives its actual moduli point.
\end{proof}
% END PUN24

\subsection{The exact remaining principal-chart boundary}

The arithmetic outputs above are complete at their printed
scope.  PUN1 supplies the literal quadratic norm theorem;
PUN10 supplies the whole original lattice model on the same
Hermitian space.  For any supplied original
\(\texttt{FamilyUnif}\), PUN11 constructs its flag-domain
models with the prescribed \(G_m\), \(\xi_0\), lattices,
principal groups, and trivial character.  PUN14 supplies
\(\texttt{LatZ}\) for every family of original data and every
cusp rank.  PUN19 supplies the original exponent match,
on the original coefficient lattice and phase, for every
\(\rho\) when \(l.A=1\).  These results do not require a
uniformization to be silently substituted for the one in the
source.

PUN20--PUN24 additionally construct actual primitive moduli
objects from every marked period, their original Weil signs
and level markings, equality precisely by marked integral
isometry, and nonemptiness at every prescribed base embedding.
They use the supplied original \(\texttt{Moduli}\) record;
existence of that representing scheme is not asserted here.
They do not yet assemble the finite family required by
\(\texttt{N1c}\).  The still-needed original producers are
\texttt{exists\_uniformisation\_at\_point\_closed},
\texttt{exists\_local\_period\_lift}, the holomorphy of the
moduli point and exponential for arbitrary holomorphic period
families, and the simultaneous theta-frame family with an
algebraic section on every fiber.  Their finite-component
assembly must preserve the same base embedding, type, level
markings, and chosen period at the base point.

After that family is constructed, the remaining original
\texttt{FramesOK}, \texttt{CoverTri}, and \texttt{FibTri}
producers must still establish one simultaneous system of
normalized framed points: the precise Gaussian norm,
triangular weight character, arithmetic-group invariance,
holomorphic readings, cover, and fiber comparison.  The
resulting \(\texttt{IsUnf}\) must use those same points,
not merely an isomorphic unframed quotient.  Finally the
\texttt{FromTest} descent must use that same family.  Only
then may the original
\texttt{FrameUnif.exists\_unf\_hasPrinc} be assembled from
these arithmetic suppliers and the analytic constructors.

Accordingly this packet supplies substantial original
subconstructors while leaving the complete principal-chart
existence statement open.  It does not close the common-field,
common-metric, common-measure, and common-phase obligations,
the full unconditional integral cusp spanning statement,
higher covers, or the all-tuple class-sum producers.
